\documentclass{article}

% ICML 2026 format (workshop papers typically use icml style)
% For submission, use the icml2026.sty if available; otherwise
% use standard article with similar margins
\usepackage{eiml_icml2026}
% \usepackage[margin=1in]{geometry}
\usepackage{amsmath,amssymb,amsfonts}
% \usepackage[T1]{fontenc}
% \usepackage{lmodern}
% \usepackage{graphicx}
\usepackage{booktabs}
% \usepackage{url}
% \usepackage{xcolor}
% \usepackage{natbib}
% \usepackage{microtype}

\title{Epistemic Blind Spots in World-Model-Based Active Inference:\\
Why RSSM Imagination Cannot See Obstacles}

\author{
Jaerock Kwon \\
Department of Electrical and Computer Engineering \\
University of Michigan--Dearborn \\
\texttt{jrkwon@umich.edu}
}

\date{}

\begin{document}
\maketitle

\begin{abstract}
We investigate whether obstacle avoidance can emerge purely from
Expected Free Energy (EFE) minimization in a deep active inference
agent for autonomous driving---without privileged obstacle
coordinates. Using an RSSM world model trained on CARLA highway
scenarios, we systematically test whether learned preferences can
drive obstacle avoidance through imagined trajectory planning. We
discover a fundamental \textbf{epistemic blind spot}: while the
RSSM posterior (filtering) correctly encodes obstacle presence,
the RSSM prior (imagination) remains obstacle-blind. A contrastive
preference model achieves 37 nats separation between obstacle-free
and obstacle-present posterior latents; an auxiliary obstacle
prediction head achieves \emph{perfect} binary separation (0.000 vs
1.000) and amplifies contrastive separation to 276 nats. Yet
obstacle avoidance remains at 0\% in all cases---the RSSM transition
model cannot predict how obstacles evolve as a function of steering
actions. This posterior-prior gap persists across three input
modalities (64$\times$64 RGB, 128$\times$128 RGB, RGBD), multiple
preference formulations, auxiliary training signals, and two data
strategies. We isolate the blind spot to the RSSM \emph{transition
model} specifically: the representation contains obstacle
information, but the dynamics model fails to propagate it through
action-conditioned rollouts. Our results demonstrate that posterior
epistemic competence does not guarantee planning-level epistemic
competence---a critical limitation for safety-relevant active
inference.
\end{abstract}

%----------------------------------------------------------------------
\section{Introduction}

Active inference \citep{friston2010free, parr2022active} provides a
unified framework for perception, learning, and action through free
energy minimization. A key promise is that goal-directed behavior
emerges from the agent's generative model without hand-crafted
reward functions: the agent imagines future trajectories, scores
them by Expected Free Energy (EFE), and selects actions that
minimize surprise relative to its preferences.

For this promise to hold in safety-critical domains like autonomous
driving, the agent's world model must be \emph{epistemically
competent}---it must know what it knows and, crucially, represent
what it doesn't know. When an obstacle appears on the road, the
world model should encode this as a deviation from preferred states,
and imagined trajectories should reflect whether actions resolve or
worsen the situation.

We test this assumption using a deep active inference agent with a
Recurrent State-Space Model (RSSM) \citep{hafner2023mastering}
world model, GMM preference distributions, and iCEM planning,
evaluated on obstacle avoidance in CARLA 0.9.16
\citep{dosovitskiy2017carla}. Our contributions:

\begin{enumerate}
    \item We identify an \textbf{epistemic blind spot} where the
    RSSM posterior correctly encodes obstacle presence but the prior
    (imagination) cannot predict obstacle dynamics from actions.

    \item We introduce a \textbf{contrastive preference model}
    (log-ratio of two GMMs) that achieves 37+ nats separation
    between obstacle-free and obstacle-present posterior latents,
    demonstrating that the representation \emph{contains} obstacle
    information.

    \item We show this posterior-prior gap is \textbf{robust across
    modalities}: 64$\times$64 RGB, 128$\times$128 RGB, and RGBD
    (depth camera) all exhibit the same failure pattern.

    \item We demonstrate that \textbf{data quality matters but is
    insufficient}: collision-free obstacle demonstrations improve
    world model driving quality (MLD 0.15m) but do not close the
    imagination gap.

    \item We \textbf{isolate the blind spot to the transition model}
    via an auxiliary obstacle prediction head that achieves perfect
    posterior separation (0.000 vs 1.000) and 276 nats contrastive
    gap---yet avoidance remains at 0\%, proving the representation
    is adequate but the dynamics model cannot propagate obstacle
    information through action-conditioned rollouts.
\end{enumerate}

%----------------------------------------------------------------------
\section{Background}

\subsection{Active Inference and EFE}

Under active inference, an agent selects policies $\pi$ that
minimize Expected Free Energy:
\begin{equation}
    G(\pi) = \sum_t \gamma^t \Big[
        \underbrace{\mathbb{E}_q[-\log p_{\text{pref}}(z_t)]}_{\text{instrumental}}
        - \underbrace{H[p(o_t | s_t)]}_{\text{epistemic}}
    \Big]
\end{equation}
where the instrumental term drives the agent toward preferred
states and the epistemic term promotes information-seeking. For
obstacle avoidance, if $p_{\text{pref}}$ assigns low probability
to obstacle-present states, the instrumental term should naturally
drive evasive actions.

\subsection{RSSM World Model}

The RSSM maintains deterministic state $h_t \in \mathbb{R}^{256}$
(GRU) and stochastic state $z_t \in \mathbb{R}^{64}$ (Gaussian):
\begin{align}
    \text{Posterior:} \quad z_t &\sim q(z_t | h_t, o_t) & \text{(filtering)} \\
    \text{Prior:} \quad \hat{z}_t &\sim p(z_t | h_t) & \text{(imagination)}
\end{align}
During \textbf{filtering} (observation available), the posterior
incorporates visual evidence---including obstacles. During
\textbf{imagination} (planning), only the prior is available,
conditioned on the deterministic state and actions. The epistemic
blind spot arises when the posterior encodes information that the
prior cannot reproduce.

%----------------------------------------------------------------------
\section{Experimental Setup}

\textbf{Architecture.} ConvEncoder (4-layer CNN, 64$\times$64
$\rightarrow$ 256D), RSSM (deter=256, stoch=64), ObsDecoder,
StateDecoder, 5-head ensemble for epistemic uncertainty. iCEM
planner (500 samples, 50 elites, 5 iterations, horizon 15).

\textbf{Tasks.} CARLA 0.9.16 highway: Task~A (lane keeping,
Town04, moderate curves) and Task~B (obstacle avoidance, Town06,
2--3 static obstacles per route, 373--518m).

\textbf{Training data.} 96K frames obstacle-free expert driving
(Town04 + Town06) plus 10--30K frames obstacle-encounter data
(BehaviorAgent and scripted lane-change demonstrations).

\textbf{Privileged baseline.} Task~B with runtime ground-truth
obstacle coordinates achieves 100\% avoidance (28/28 obstacles)
via reflexive steer prior, motor primitives, and directional CEM
penalty. All six privileged mechanisms are removed for the pure
AIF experiments.

%----------------------------------------------------------------------
\section{Results}

\subsection{The Preference Gap is Inverted}

We fit a GMM preference model on obstacle-free expert latents and
evaluate log-probability on held-out obstacle-free vs.\
obstacle-present latents. If the GMM correctly encodes ``prefer
obstacle-free states,'' obstacle-present latents should have lower
log-probability.

\begin{table}[h]
\centering
\caption{Preference gap: $\Delta = \log p(z_{\text{clean}}) - \log p(z_{\text{obs}})$.
Positive = correct direction.}
\label{tab:gap}
\begin{tabular}{lcc}
\toprule
\textbf{Input modality} & \textbf{$\Delta$ (nats)} & \textbf{Recon ratio} \\
\midrule
64$\times$64 RGB & $-$24.6 & 1.62$\times$ \\
128$\times$128 RGB & $-$2.3 & 1.91$\times$ \\
64$\times$64 RGBD & $-$68.1 & 1.35$\times$ \\
\bottomrule
\end{tabular}
\end{table}

The gap is \textbf{inverted in all cases}: obstacle-present latents
receive \emph{higher} GMM probability. The RSSM latent space
encodes driving dynamics (speed, steering, lane position) rather
than scene content. Higher resolution (128$\times$128) narrows the
gap but does not reverse it. Adding depth worsens it.

\subsection{Contrastive Preference Achieves Posterior Separation}

We replace the single GMM with a contrastive model:
$\log p_{\text{pref}}(z) = \log p_{\text{clean}}(z) - \alpha \cdot \log p_{\text{avoid}}(z)$,
where $p_{\text{clean}}$ is fitted on obstacle-free and
$p_{\text{avoid}}$ on obstacle-present latents.

\begin{table}[h]
\centering
\caption{Contrastive preference: posterior separation is correct,
but obstacle avoidance remains at 0\%.}
\label{tab:contrastive}
\begin{tabular}{lccc}
\toprule
\textbf{Scale $\alpha$} & \textbf{Raw gap} & \textbf{MLD (m)} & \textbf{Avoidance} \\
\midrule
1.0 & +37.0 nats & 4.0 & 0\% (off-road) \\
0.3 & +9.2 nats & 0.32 & 0\% \\
0.2 & +1.1 nats & 0.17 & 0\% \\
0.1 & +1.0 nats & 0.31 & 0\% \\
\bottomrule
\end{tabular}
\end{table}

At $\alpha{=}1.0$, the 37 nats gap is strong but overwhelms
driving (MLD 4.0m). At $\alpha{\leq}0.3$, driving is preserved
but avoidance remains 0\%. The contrastive signal correctly
separates \emph{posterior} latents but cannot influence
\emph{imagined} trajectories.

\subsection{The Epistemic Blind Spot}

The root cause: during CEM planning, the RSSM imagines 500
trajectories via \texttt{img\_step} (prior only). These imagined
latent sequences do not encode obstacle proximity because the
prior cannot predict obstacle appearance/disappearance from
steering actions. All trajectories---straight, left, right---produce
similar latent sequences in imagination.

\textbf{Evidence:} The 128$\times$128 RGB model achieves 100\%
success rate (goal reached) with 0.17m mean lateral deviation and
zero offroad events---the best driving quality of any variant---yet
avoids zero obstacles. The agent drives perfectly but is blind to
obstacles in its plans.

\subsection{Clean Data Does Not Close the Gap}

We collected 30K collision-free obstacle avoidance demonstrations
(40\% clean rate via waypoint-redirect autopilot) and retrained.
With 10K obstacle frames (9.4\% of training data), driving quality
is preserved (MLD 0.15m) but avoidance remains 0\%. The RSSM
learns to \emph{reconstruct} obstacle images (1.62$\times$ higher
reconstruction error) but not to \emph{imagine} obstacle dynamics.

\subsection{Auxiliary Head Isolates the Transition Model}

To determine whether the blind spot lies in the \emph{representation}
(latent space lacks obstacle information) or the \emph{transition
model} (dynamics cannot propagate it), we add an auxiliary obstacle
prediction head---a small MLP (320$\rightarrow$128$\rightarrow$1)
trained with binary cross-entropy alongside the VFE loss. The head
receives RSSM features and predicts obstacle presence
($\beta_{\text{aux}}{=}1.0$). Crucially, this head is used
\textbf{only during training} to shape the latent space; at
evaluation time, only the shaped representation and standard EFE
scoring are used.

\begin{table}[h]
\centering
\caption{Auxiliary head forces obstacle information into the latent
space, achieving perfect posterior separation. Avoidance remains 0\%.}
\label{tab:aux}
\begin{tabular}{lccc}
\toprule
\textbf{Model} & \textbf{Obs.\ head} & \textbf{Contr.\ gap} & \textbf{Avoid} \\
 & (clean / obs) & (eff.\ nats) & \\
\midrule
v6 (no aux) & N/A & 1.1 & 0\% \\
v10 (no aux, clean data) & N/A & 1.1 & 0\% \\
v11 (aux head) & 0.000 / 1.000 & 133.1 & 0\% \\
\bottomrule
\end{tabular}
\vspace{1mm}
\small{Obs.\ head: sigmoid output (0=clean, 1=obstacle).
Contr.\ gap: contrastive log-ratio at $\alpha{=}0.2$.}
\end{table}

The auxiliary head achieves \textbf{perfect binary separation}
(Table~\ref{tab:aux}): 0.000 for obstacle-free frames, 1.000 for
obstacle-present frames. The contrastive preference gap increases
from 1.1 to \textbf{133.1 effective nats} (276 raw)---a
$120\times$ improvement. This proves the latent space \emph{can}
encode obstacle presence when forced by supervision.

Yet obstacle avoidance remains at \textbf{0\%}. The agent with
the auxiliary-head-shaped latent space still collides with every
obstacle. This result \emph{definitively isolates the blind spot}:

\begin{itemize}
    \item \textbf{Representation}: adequate. The latent space
    encodes obstacle presence with perfect fidelity.
    \item \textbf{Transition model}: deficient. The GRU-based
    dynamics $f(h_t, z_t, a_t) \rightarrow h_{t+1}$ does not
    learn that ``steer left when obstacle ahead'' produces
    different future latent states than ``continue straight.''
\end{itemize}

During imagination, the RSSM transition model rolls out future
states using only the prior. Even though the current posterior
state encodes ``obstacle present,'' the transition model maps all
action sequences to similar future states---it has not learned
obstacle-conditioned dynamics.

%----------------------------------------------------------------------
\section{Discussion}

\subsection{Posterior vs.\ Prior Epistemic Competence}

Our results reveal a dissociation between two levels of epistemic
competence in world-model-based active inference:

\textbf{Posterior competence} (filtering): The RSSM posterior
correctly encodes obstacle presence when conditioned on
observations. The contrastive preference achieves 37+ nats
separation without auxiliary supervision, and \emph{276 nats} with
an auxiliary prediction head. The representation is not the
bottleneck.

\textbf{Prior incompetence} (imagination): The RSSM prior cannot
predict how obstacle-related visual features evolve as a function
of steering actions. Even when the auxiliary head forces perfect
obstacle encoding into the latent space (0.000 vs 1.000), the
transition model does not propagate this information through
action-conditioned rollouts.

This posterior-prior gap is an \textbf{epistemic blind spot}
localized in the \textbf{transition model}: the agent can
recognize danger when it sees it but cannot predict danger from
its plans. In active inference terms, the generative model's
\emph{likelihood} (observation model) encodes obstacles, but the
\emph{transition} model $f(h_t, z_t, a_t) \rightarrow h_{t+1}$
does not learn obstacle-conditioned dynamics. The auxiliary head
experiment serves as a controlled ablation that separates
representation adequacy from transition adequacy.

\subsection{Implications for AIF Safety}

This finding has direct implications for deploying active inference
agents in safety-critical settings:

\begin{enumerate}
    \item \textbf{EFE planning is only as good as imagination.}
    A preference model that perfectly separates safe from unsafe
    states is useless if the world model cannot imagine the
    consequences of actions in those states.

    \item \textbf{Observation-level epistemic uncertainty is
    insufficient.} Ensemble disagreement and reconstruction error
    detect obstacles at the observation level, but these signals
    are not available during imagination (no observations).

    \item \textbf{The latent bottleneck prioritizes dynamics over
    content.} With 64 stochastic dimensions, the RSSM encodes
    temporally smooth driving dynamics rather than spatially
    localized scene content (obstacles). This is efficient for
    lane keeping but blind to safety-critical objects.
\end{enumerate}

\subsection{Toward Closing the Gap}

Three directions may address this epistemic blind spot while
remaining within the active inference framework:

\textbf{Larger latent space.} Increasing stochastic dimensions
from 64 to 256 may provide enough capacity to encode both driving
dynamics and scene content. This is the most self-supervised
approach but may require proportionally more training data.

\textbf{Auxiliary prediction signal.} We tested a supervised
obstacle detection head during training that forces the latent
space to encode obstacle presence (Section~4.5). While this
achieves perfect posterior separation, the transition model
remains obstacle-blind. A stronger version---training the
auxiliary head on \emph{imagined} rollouts rather than posterior
states---may force the transition model to propagate obstacle
information, but this requires differentiating through the
imagination loop.

\textbf{Object-centric architectures.} Models like SLATE
\citep{singh2022illiterate} or SAVi \citep{kipf2022conditional}
that decompose scenes into object slots may naturally separate
obstacles from background, making obstacle dynamics explicitly
representable in imagination.

%----------------------------------------------------------------------
\section{Related Work}

\textbf{Deep active inference.} \citet{catal2020learning} and
\citet{fountas2020deep} apply active inference with VAE world
models to visual navigation and Atari. \citet{millidge2020deep}
connects active inference with policy gradients.
\citet{lanillos2021active} surveys active inference in robotics.
Our work extends to continuous-action driving and identifies a
specific failure mode in RSSM-based imagination.

\textbf{World models for driving.} DreamerV3
\citep{hafner2023mastering} uses RSSM imagination for RL-based
policy learning. MILE \citep{hu2022model} learns driving world
models for imitation. Our finding---that RSSM imagination is
obstacle-blind---may apply to these systems as well.

\textbf{Epistemic uncertainty in planning.}
\citet{chua2018deep} use ensemble disagreement for model-based RL.
\citet{sekar2020planning} explore planning to explore. Our work
shows that epistemic signals available during filtering (posterior)
may not transfer to planning (prior).

%----------------------------------------------------------------------
\section{Conclusion}

We identified a fundamental epistemic blind spot in RSSM-based
active inference and isolated it to the \emph{transition model}.
The world model's posterior correctly encodes obstacle
presence---a contrastive preference achieves 37 nats separation,
and an auxiliary prediction head achieves perfect binary
classification (276 nats)---but the prior (imagination) cannot
predict obstacle dynamics from actions. This posterior-prior gap
persists across input modalities (RGB, high-resolution RGB, RGBD),
preference formulations (single GMM, contrastive), auxiliary
supervision, and training data quality. The result is an agent
that drives excellently (100\% SR, 0.17m MLD at 128$\times$128)
but is completely blind to obstacles in its planning horizon.

The auxiliary head experiment provides the critical ablation:
forcing obstacle information into the latent space with perfect
fidelity does not help, because the GRU-based transition model
$f(h_t, z_t, a_t) \rightarrow h_{t+1}$ does not learn
obstacle-conditioned dynamics. This isolates the problem from
representation capacity to transition model expressiveness.

Our findings highlight that posterior epistemic competence does
not guarantee planning-level epistemic competence---a distinction
that is critical for deploying active inference agents in
safety-relevant domains. Future work should focus on transition
model architectures that can learn object-conditioned dynamics,
such as attention-based transitions or object-centric world
models, to close this posterior-prior gap.

%----------------------------------------------------------------------
\bibliographystyle{plainnat}
\bibliography{references}

\end{document}
